\underline{定理} 2.5. \underline{沿用 2.1 的记号:}

\begin{enumerate}[label=(\roman*),leftmargin=2.6em]
\item \underline{若 $p+q=p'+q'$ 且 $p\leq p'\leq q'\leq q$, 则}
\begin{equation*}
h_0^{p,q}(\underline a)\leq h_0^{p',q'}(\underline a).
\end{equation*}

\item \underline{若 $p\leq q$, 则 $h_0^{p,q}(\underline a)\ne0$ 当且仅当}
\begin{equation*}
p>\left\lfloor
\frac{p+q+d-\sum_i a_i}{\sup_i(a_i)}
\right\rfloor.
\end{equation*}
\underline{这里 $\lfloor\ \rfloor$ 表示商的整数部分.}
\end{enumerate}

%%%%%%%%%% source scan idx 62 | folio 55 %%%%%%%%%%

设 (*) 表示形式幂级数 $P=\sum a_{p,q}y^pz^q$ 的下列性质:

\begin{quote}
(*) $a_{p,q}$ 非负, $a_{p,q}=a_{q,p}$, 且只要 $p+q=p'+q'$ 及 $p\leq p'\leq q'\leq q$, 就有 $a_{p,q}\leq a_{p',q'}$.
\end{quote}

由 2.4 中的公式、恒等式 $(1-x)^{-1}=\sum_{n\geq0}x^n$ 以及下面的引理, 立即得到断言 (i). 引理的证明留给读者.

\underline{引理} 2.5.1. \underline{满足 (*) 的形式幂级数的集合对于加法和乘法稳定.}

先在超曲面 ($d=1$) 这一特殊情形下证明 (ii). 此时对于 $p\leq q$, 欲证公式可改写为

\begin{equation}\tag{2.5.2}
h_0^{p,q}(a)\ne0
\quad\Longleftrightarrow\quad
p\geq\left\lfloor\frac{p+q+1}{a}\right\rfloor
\end{equation}

或等价地, $(p+1)a\geq p+q+2$. 由 2.4(i) 中的第二个公式可见, $h_0^{p,q}(a)$ 非零当且仅当存在数对 $(e,f)$ 和 $(c_i,d_i)$, 使得

\begin{equation}\tag{2.5.3}
\left\{
\begin{aligned}
(e,f)&\geq(0,0),& e+f&\leq a-2,\\
(c_i,d_i)&\geq(1,1),& c_i+d_i&\leq a,\\
(p,q)&=(e,f)+\sum_i(c_i,d_i).&&
\end{aligned}
\right.
\end{equation}

%%%%%%%%%% source scan idx 63 | folio 56 %%%%%%%%%%

固定 $p+q$ 时, 取 $e=0$ 且各 $c_i=1$, 就得到尽可能小的 $p$; 此值就是指标 $i$ 的个数. 因而它是满足下式的 $p$ 的下界

\begin{equation*}
p\geq\frac{(p+q)-(a-2)}{a},
\qquad\text{即}\qquad (p+1)a\geq p+q+2.
\end{equation*}

现在转入一般情形. 若 $\underline a'$ 是从 $\underline a$ 中删去等于一的 $a_i$ 所得, 则 $H(\underline a)=H(\underline a')$. 从 2.5(ii) 的形式看, 因此可假设每个 $a_i>1$. 若 $d=0$, 则 $h_0^{p,q}=0$, 断言成立; 故不失一般性, 可假设 $d>0$ 且 $a_d=\sup_i(a_i)$.

由 2.4(ii) 和 (2.5.2), 条件 $h_0^{p,q}(\underline a)\ne0$ 等价于存在 $[1,d]$ 的非空子集 $P$、数对 $(e,f)$ 以及数对 $(c_i,d_i)_{i\in P}$, 使得

\begin{equation}\tag{2.5.4}
\left\{
\begin{aligned}
(|P|-1,|P|-1)&\geq(e,f)\geq(0,0),\\
c_i,d_i&\geq
\left\lfloor\frac{c_i+d_i+1}{a_i}\right\rfloor,
& (c_i,d_i)&\geq(0,0),\\
(p,q)&=(e,f)+\sum_{i\in P}(c_i,d_i).&&
\end{aligned}
\right.
\end{equation}

固定 $n=p+q$, 求 $p$ 的最小可能值. 置 $n_i=c_i+d_i$. 对于这样的 $p$, 条件 (2.5.4) 变为

%%%%%%%%%% source scan idx 64 | folio 57 %%%%%%%%%%

\begin{equation}\tag{2.5.5}
\left\{
\begin{aligned}
a_i=2&\Longrightarrow n_i\text{ 为偶数}
&&\text{(否则 (2.5.4) 不可能成立)},\\
c_i&\geq\left\lfloor\frac{n_i+1}{a_i}\right\rfloor,
&c_i&\leq n_i,\\
(|P|-1,|P|-1)&\geq(e,f)\geq(0,0),&&\\
n&=e+f+\sum_{i\in P}n_i,&&\\
p&=e+\sum_{i\in P}c_i.&&
\end{aligned}
\right.
\end{equation}

由此描述首先可见, 最小的 $p$ 在 $P=[1,d]$ 时取得, 以下即限于此情形. 若

\begin{equation}\tag{2.5.6}
\sum_i(a_i-2)+(d-1)\geq n,
\end{equation}

则存在满足 $p=c_i=e=0$ 的解. 否则, 若 $a_d\ne2$, 最佳选择是对 $i\ne d$ 取 $n_i=a_i-2$, 因而当 $i\ne d$ 时 $c_i=0$, 且 $e=0$, $f=d-1$, 以及

\begin{equation*}
n_d=n-(d-1)-\sum_{i\ne d}(a_i-2).
\end{equation*}

于是

\begin{equation}\tag{2.5.7}
p=\left\lfloor
\frac{n-(d-1)-\sum_{i\ne d}(a_i-2)+1}{a_d}
\right\rfloor,
\end{equation}

所以, 对于 $p\leq q$, 有 $h_0^{p,q}(\underline a)\ne0$ 当且仅当

\begin{equation}\tag{2.5.8}
p\geq\left\lfloor
\frac{n-\sum_{i\ne d}(a_i-1)+1}{a_d}
\right\rfloor.
\end{equation}

当 $d=1$ 时, 公式 (2.5.8) 化为 (2.5.2), 而当 (2.5.6) 成立时它显然成立. 若 $d\geq2$, (2.5.6) 不成立, 且所有 $a_i$ 均等于二, 则最佳选择改为

%%%%%%%%%% source scan idx 65 | folio 58 %%%%%%%%%%

当 $i\ne d$ 时取 $n_i=0$, 取 $e=0$, 并取 $f=d-1$ 或 $d-2$, 使其与 $n$ 模二同余, 再令 $n_d=n-f$. 若 $n\equiv d-1\pmod2$, 则与上述相同可证 (2.5.7). 否则, 注意到

\begin{equation*}
\left\lfloor\frac{n-(d-2)+1}{2}\right\rfloor
=\left\lfloor\frac{n-(d-1)+1}{2}\right\rfloor,
\end{equation*}

故当 $a_d=2$ 时, (2.5.7) 从而 (2.5.8) 仍成立.

只须再验证 (2.5.8) 与 2.5(ii) 等价. 事实上,

\begin{equation*}
\left\lfloor\frac{n-\sum_{i\ne d}(a_i-1)+1}{a_d}\right\rfloor
=\left\lfloor\frac{n+d-\sum_{i\ne d}a_i}{a_d}\right\rfloor
=\left\lfloor\frac{p+q+d-\sum_i a_i}{a_d}\right\rfloor+1.
\end{equation*}

2.6. 若 $X$ 是 $\mathbb P^r$ 中维数为 $n$ 的完全交, 则由 (2.5.8) 立即可知, 只有当 $X$ 为奇数维二次超曲面时, $H^n(X)$ 的本原部分才为零.

